\section{Introduction}

Vol.~1 of this series established a minimal framework for the emergence of
time and space from the causal log $C_{ij}$, defined on a network of nodes
subject to the single axiom \emph{there is difference}~\cite{VED_vol1}.
The symmetric and antisymmetric decomposition
\begin{equation}
  S_{ij} = \tfrac{1}{2}(C_{ij}+C_{ji}), \qquad
  A_{ij} = \tfrac{1}{2}(C_{ij}-C_{ji})
\end{equation}
encodes spatial and causal structure within the same variable.

The present volume addresses the next question:
given that space and time emerge from closure geometry,
can the internal structure of matter --- particles, symmetries, charges, and masses ---
be understood within the same framework?

The Standard Model of particle physics describes this internal structure
with remarkable precision.
Its gauge group $\mathrm{SU}(3)\times\mathrm{SU}(2)\times\mathrm{U}(1)$,
particle spectrum, charge assignments, generation structure,
and mass-generation mechanism are experimentally confirmed,
yet their geometric or structural origin remains unexplained.

The approach taken here is the following.
Particles are interpreted not as fundamental entities,
but as stable configurations of closure structure.
Symmetries are identified not as a priori principles,
but as transformation groups that leave closure invariant.
Charges, generations, and mass arise as features of the internal geometry
of triangular closure.

The central claim of this volume is:
\begin{equation}
  \boxed{
    \text{The Standard Model can be interpreted as a structural consequence
    of closure geometry.}
  }
\end{equation}

This interpretation is structural and qualitative in nature.
It does not constitute a predictive theory,
and no claims are made about numerical coincidences with experimental values.
The goal is to exhibit a coherent geometric reading of the Standard Model's skeleton.

The paper is organized as follows.
Section~\ref{sec:closure_energy} introduces the closure energy functional
and identifies particles as stable configurations.
Section~\ref{sec:triangular} establishes the triangle as the unique
minimal closure unit.
Section~\ref{sec:symmetry} derives the gauge symmetry structure
from closure invariance.
Sections~\ref{sec:particles}--\ref{sec:higgs} treat
particle classification, generation structure, electric charge,
the weak interaction, and the Higgs mechanism in turn.
Section~\ref{sec:discussion} discusses the relation to Vol.~1
and the path toward Vol.~3.
